\documentclass[10pt,a4paper]{amsart}
\usepackage{amsmath,amssymb,amsthm,graphicx}
\RequirePackage{amssymb,amsmath,amsthm,bm}
\RequirePackage{thmtools}
\RequirePackage[capitalize]{cleveref}
\providecommand{\tightlist}{%
  \setlength{\itemsep}{0pt}\setlength{\parskip}{0pt}}

%%%%% varphi instead of phi
\renewcommand{\phi}{\varphi}

\mathchardef\mhyph="2D 
\providecommand{\A}{\mathcal{A}}
\providecommand{\C}{\mathcal{C}}
\providecommand{\L}{\mathcal{L}}
\providecommand{\B}{\mathcal{B}}
\providecommand{\M}{\mathcal{M}}
\renewcommand{\S}{\mathcal{S}}
\providecommand{\ol}[1]{\overline{#1}}
\providecommand{\mbf}[1]{\mathbf{#1}}
\providecommand{\mf}[1]{\mathfrak{#1}}
\providecommand{\mrm}[1]{\mathrm{#1}}
\providecommand{\mc}[1]{\mathcal{#1}}
\providecommand{\ra}{\rightarrow}
\providecommand{\Ra}{\Rightarrow}
\providecommand{\La}{\Leftarrow}
\providecommand{\lr}{\leftrightarrow}
\providecommand{\hf}{\mathbb{HF}}
\providecommand{\spec}{DgSp}
\providecommand{\bSigma}{\bm\Sigma}
\providecommand{\bDelta}{\bm\Delta}
\providecommand{\bPi}{\bm\Pi}
\providecommand{\Pinf}[1]{\Pi^{\mathrm{in}}_{#1}}
\providecommand{\Sinf}[1]{\Sigma^{\mathrm{in}}_{#1}}
\providecommand{\dSinf}[1]{d\mhyph\Sigma^{\mathrm{in}}_{#1}}
\providecommand{\Dinf}[1]{\Delta^{\mathrm{in}}_{#1}}
\providecommand{\bp}{\bar p}
\providecommand{\bq}{\bar q}
\providecommand{\ba}{\bar a}
\providecommand{\bb}{\bar b}
\providecommand{\bc}{\bar c}
\providecommand{\bd}{\bar d}
\providecommand{\bx}{\bar x}
\providecommand{\by}{\bar y}
\providecommand{\bz}{\bar z}
\providecommand{\LO}{\mathrm{LO}}
\declaretheorem[name={Theorem},numberwithin=section]{theorem}
\declaretheorem[name={Proposition}, sibling=theorem]{proposition}
\declaretheorem[name={Lemma}, sibling=theorem]{lemma}
\declaretheorem[name={Fact}, sibling=theorem]{fact}
\declaretheorem[name={Claim}, numberwithin=theorem]{claim}
\declaretheorem[name={Definition},numberwithin=section,sibling=theorem, style=definition]{definition}
\declaretheorem[name={Remark}, numbered=no,style=definition]{remark}
\declaretheorem[name={Corollary}, sibling=theorem]{corollary}
\declaretheorem[name={Question},sibling=theorem]{question}
\newcommand{\T}{{\mathcal{T}}}
\newcommand{\SR}{SR}
\newcommand{\pSR}{pSR}
\newcommand{\Q}{\mathbb{Q}}
\newcommand{\Z}{\mathbb{Z}}
\newcommand{\K}{\mathcal{K}}
\newcommand{\N}{\mathcal{N}}
\renewcommand{\L}{\mathcal{L}}
\providecommand{\Ginf}[1]{\Gamma^{\mathrm{in}}_{#1}}
\title{Inserting an order into a shuffle sum lowers its back-and-forth rank by two}
\author{David Gonzalez, Dino Rossegger}
\date{arxiv:2305.07126v1; CC BY 4.0}
\begin{document}
\maketitle
\noindent\textit{Editorial note.} The counted statement is the author environment at source0.tex lines 1024-1027 and its proof environment opens at line 1028, the line immediately after; the pairing is by adjacency and the retained proof range 1028-1056 is the whole environment.\par\medskip
\noindent\textit{Editorial note.} Context: the author introduces the shuffle sum of a set of countable linear orderings, and remarks that an order of the form L times Q is a special case, at source0.tex lines 1016-1023; that passage is retained as a context span because the counted statement and its proof use both the construction and that special case. It is the author wording, retained verbatim, and it is not counted.\par\medskip
\noindent\textit{Editorial note.} The retained preamble is the author style file dino.sty (package lines, the theorem declaration support, the notation and tuple macros and the theorem environments) followed by the ordinal and rank macros of the manuscript itself, reproduced verbatim; this is the same recipe as the previously published card of this article, so no package outside it is requested and extra\_packages stays empty.\par\medskip
\noindent\textit{Editorial note.} Metadata provenance: source\_id and source\_version are the values recorded in this source licence evidence file (sources/scott-linear-orders-2305.07126v1/LICENSE\_EVIDENCE.json: arxiv:2305.07126v1 / 2305.07126v1), which the card points at directly; the title and author list are transcribed from the authors\_index and title keys of the source record of the pinned copy. The card carries no counter setting and imposes no original numbering, since this article has no label-number dump and its numbering cannot be reconstructed from the retained ranges alone.\par\medskip
\noindent\textit{Editorial note.} No external reference is needed: apart from the definition retained as context, the statement and proof use no label or citation from outside their own ranges. The manuscript continues after the retained proof with a second lemma that repeats this statement under an additional hypothesis; that later environment is outside the retained range and is not counted.\par\medskip
\noindent\textit{Editorial scope.} One original result of the pinned author manuscript, counted once: the rank alpha plus two comparison obtained by inserting an order into a shuffle sum (author label lem:shufflesum).\par\medskip
\noindent\textit{Allowed standard prerequisites (editorial).}\begin{itemize}
\item Shuffle sum. The shuffle sum of a family of countable linear orders is obtained by partitioning the rationals into mutually dense sets and replacing each point of one of them by a copy of the corresponding order.. It is the construction to which the counted order belongs, and its definition is retained in the card.
\item Order type L times Q. An order of the form L times Q is a shuffle sum, and every interval of it is again of that form.. It is the shape of the second order in the statement.
\item Splitting a shuffle sum at a point. Above any point of a shuffle sum the remainder is again a shuffle sum, and on the other side the order is again of shuffle sum form; in particular a shuffle sum is isomorphic to itself with a copy of itself appended.. It is what allows the two reductions in the second half of the proof.
\end{itemize}\medskip
our examples are constructed by producing shuffle sums of linear orderings.
Recall that given a set $S$ of countable linear orderings, the \emph{shuffle
sum} $Sh(S)$ of $S$ is obtained by partitioning $\Q$ into $|S|$ mutually dense
sets $K_i$ and replacing each point in $K_i$ with a copy of $S_i$.
Note that a linear order of the form $L\cdot\Q$ is a special case of a shuffle
sum. We begin by proving a couple of lemmas about the back-and-forth relations
of shuffle sums.


\begin{lemma}\label{lem:shufflesum}
Let $\A$ be a shuffle sum and $\B = \L\cdot\Q$ for some ordering $\L$. If
$\A\leq_\alpha \B$ then $\A+\L+\B\leq_{\alpha+2}\A+\B$.
\end{lemma}

\begin{proof}
To show this, we describe a winning strategy for the $\exists$-player in the
$\alpha+2$-back-and-forth game. 
On the first turn, the $\exists$-player copies the selected tuple of the
$\forall$-player using an isomorphism of the initial $\A$s and the final $\B$s. 
The only possible non-isomorphic interval is the one between the smallest
element $b$ in $\B$ and the largest element $a$ in $\A$. 
In other words, we need to show that 
\[\A_{>a}+\L+\B_{<b}\geq_{\alpha+1} \A_{>a}+\B_{<b}.\]
Because $\A$ and $\B$ are shuffle sums, $\A_{>a}\cong N+\A$ and $\B_{<b}\cong
\B+K$ for some linear orderings $N$ and $K$. 
This means that we only need to show that
$\A+\L+\B\geq_{\alpha+1}\A+\B$
if we play isomorphically on $N$ and $K$.

On the next turn, the $\forall$-player plays on $\A+\L+\B$, and the
$\exists$-player responds on $\A+\B$ using an isomorphism between the initial
$\A$s and an isomorphism of $\L+\B$ to a final segment of $\B$.
Assume without loss of generality, that $c$ is the smallest element in $\L$
played by the $\forall$-player in the distinguished copy of $\L$ within $\A+\L+\B$.
The only possible non-isomorphic interval is between $c$ and the
largest element $a$ selected within $\A$.
In other words, we need to show that 
\[\A_{>a}+\L_{<c}\leq_{\alpha} \A_{>a}+\B+\L_{<c}.\]
As $\A_{>a}\cong M+\A$ and the $\exists$-player can play the isomorphism between
the copies of $\L_{<c}$ and $M$, it is sufficient to show that $\A\leq_{\alpha}\A+\B$.
As $\A$ is a shuffle sum $\A\cong \A+\A$.
Therefore, it is enough to show that $\A\leq_\alpha \B$, as desired.
\end{proof}

\end{document}
